\documentclass[11pt]{article}

\usepackage[margin=1in]{geometry}
\usepackage[T1]{fontenc}
\usepackage[utf8]{inputenc}
\usepackage{setspace}
\usepackage{hyperref}

\onehalfspacing

\begin{document}

\begin{center}
{\Large \textbf{Adaptive Wildfire Governance Through Reputation-Weighted Response Networks}}\\[1em]

{\large Fabian Leo Naressi$^{*}$}\\
{\large Independent researcher}\\[0.5em]
{\small $^{*}$Correspondence: \href{mailto:fabian.leo.naressi@example.com}{fabian.leo.naressi@example.com}}
\end{center}

\vspace{1.5em}

\noindent \textbf{Keywords:} wildfire governance; reputation systems; emergency response; resilience; adaptive coordination\\
\textbf{JEL classification:} Q54; Q58

\vspace{1.5em}

\begin{abstract}
This paper proposes an adaptive governance framework for wildfire response based on reputation-weighted coordination among distributed response actors. The model combines reputation scoring, resource allocation, and a circulation-inspired incentive mechanism to improve the timeliness, resilience, and accountability of wildfire governance in high-risk landscapes. A lightweight simulation of representative fire-management agencies demonstrates how reputation-weighted decision influence can support faster consensus on resource deployment and more transparent post-event evaluation. The paper contributes a reproducible conceptual prototype that bridges wildfire governance, trust-weighted collaboration, and adaptive response design.
\end{abstract}

\end{document}
